\documentclass[10pt,a4paper]{amsart}
\usepackage[margin=19mm]{geometry}
\usepackage{amsmath,amssymb,mathtools}
\usepackage[scr=rsfs,cal=euler]{mathalfa}
\usepackage{xfrac}
\usepackage[unicode,hidelinks,bookmarks=false]{hyperref}
\newcommand{\ie}{i.e.\ }
\newcommand{\cf}{cf.\ }
\newcommand{\resp}{respectively\ }
\newcommand{\Subsection}{\subsection}
\providecommand{\nobreakdash}{\nobreak\hbox{-}}
\setlength{\emergencystretch}{2em}
\multlinegap0pt
\usepackage{amssymb}
\usepackage{bm}
\usepackage{xcolor}
\usepackage{algorithm}\usepackage{algpseudocode}
\usepackage{tikz}
\newtheorem{lemma}{Lemma}
\newtheorem{proposition}{Proposition}
\newcommand{\bbE}{\mathbb{E}}
\newcommand{\calE}{\mathcal{E}}
\newcommand{\calN}{\mathcal{N}}
\newcommand{\hatb}{\hat{b}}
\newcommand{\nn}{\nonumber} 
\newcommand{\rvC}{\mathsf{C}}
\title{Second-Order Asymptotics for the Gaussian Multiple-Access Channel at Corner Points}
\author{Vincent Y. F. Tan}
\date{Source: arXiv:2608.18231v1 (CC BY 4.0)}
\begin{document}
\maketitle

\noindent\emph{Source and licence.} Second-Order Asymptotics for the Gaussian Multiple-Access Channel at Corner Points. Version arXiv:2608.18231v1 (CC BY 4.0), distributed by arXiv under the Creative Commons Attribution 4.0 International licence, http://creativecommons.org/licenses/by/4.0/, as recorded in the arXiv licence metadata for this exact version and shown on the abstract page https://arxiv.org/abs/2608.18231v1. Full licence notice: \texttt{licenses/arxiv-2608.18231-CC-BY-4.0.txt}.\par\smallskip

\noindent\textit{Editorial scope.} Counted unit: the article's proposition prop:excep\_likeli (source0.tex lines 2549--2555). The article's complete original proof of it is delivered with it (source0.tex lines 2558--2622, one \texttt{proof} environment of 65 lines). No mathematics is written, repaired or re-derived: the statement and the proof are the article's own lines, in the article's own notation, and no retained line is substituted or re-flowed.  Retained uncounted as the source-local context the proof needs: lines 2523--2537.  Nothing was omitted from any retained span, so the package carries no omission record.  No retained line contains a citation, so the wrapper carries no bibliography and no bibliography entry.  No retained line contains a figure environment, an image-inclusion command or drawing code, so no image is delivered and the package declares no asset dependency.  Editorial formatting only, in addition: an AMS \texttt{amsart} preamble that restates the article's own declarations and macros used by the retained lines, namely the environments \texttt{lemma}, \texttt{proposition} on separate counters, together with the standard packages those lines need.  The retained environments print in this document as Lemma 1, Proposition 1 in order of appearance, whereas the article prints the same items with the numbering of its own full document.  The complete native source of the article as distributed in the arXiv e-print for this exact version is bundled unchanged as \texttt{sources/207976/source0.tex}.
\begin{lemma}[Broad Gaussian Likelihood] \label{lem:broad_gauss}
    Let $\calE_n$ have dimension $r\ge1$, let $B\ge 0$ and suppose $x^r\in \calE_n$ satisfies $\|x^r\|^2\le B$. Let $Z^r\sim\calN(0,I_{\calE_n})$, $Y^r=x^r+Z^r$ and 
    \begin{align}
        q_B (y^r)= \calN\left(y^r;0,\Big( 1+\frac{B}{r}\Big) I_{\calE_n}\right).
    \end{align}
    Then, the log-likelihood ratio 
    \begin{align}
        \Lambda_B(x^r,Z^r) &= \log\frac{\calN(Y^r;x^r, I_{\calE_n})}{q_B  (Y^r)} \nn\\
        &=r\rvC(B/r)+ \frac{ \|x^r\|^2-B}{2(1+B/r)}+ \frac{\langle x^r,Z^r\rangle}{1+B/r}- \frac{B/r}{2(1+B/r)} \big(\|Z^r\|^2-r\big) \label{eqn:Lambda_B}
    \end{align}
    satisfies
    \begin{align}
        \bbE[ \Lambda_B(x^r,Z^r)]\le r\rvC(B/r)\quad\mbox{and}\quad \mathrm{Var}( \Lambda_B(x^r,Z^r))\le\frac{r}{2}.
    \end{align}
\end{lemma}

\begin{proposition} \label{prop:excep_likeli}
    For $\hatb_n \in \{b_n,4b_n\}$, let $Z_\calE^n\sim\calN(0,I_{\calE_n})$. With $\Lambda_{\hatb_n}$ defined in Lemma~\ref{lem:broad_gauss}  and interpreted as zero when $r_n=0$, for every $\eta>0$,
    \begin{align}
        \sup_{x^n \in \calE_n : \|x^n\|^2\le\hatb_n}\Pr\big( \Lambda_{\hatb_n}(x^n,Z_\calE^n)^+>\eta\sqrt{n} \big)\to0. \label{eqn:prop_cheby}
    \end{align}
    Here and in the following, $u^+:=\max\{u,0\}$. 
\end{proposition}

\begin{proof}[Proof of Proposition~\ref{prop:excep_likeli}]
Recall that $b_n=n^{7/8}$. If $r_n=0$, then $\calE_n=\{0\}$ and the
exceptional likelihood is defined to be zero. It therefore suffices to
consider indices for which $r_n\ge1$.

For $B>0$, define $f_B(r):=r\rvC(B/r)$ for $ r>0$.
This function is nondecreasing because
\begin{align}
f_B'(r)
=\frac12\left[
\log\left(1+\frac Br\right)-\frac{B}{B+r}
\right]\ge0.
\end{align}
Choose a finite constant $K_0$ such that
$r_n\le K_0n^{3/8}$ for all sufficiently large $n$, and put $R_n:=K_0n^{3/8}$.
For either choice $B\in\{b_n,4b_n\}$, monotonicity gives
\begin{align}
\Upsilon_n(B)
&:=r_n\rvC(B/r_n)\le R_n\rvC(B/R_n)
\le
\frac{K_0n^{3/8}}{2}
\log\left(1+\frac{4b_n}{K_0n^{3/8}}\right)=O(n^{3/8}\log n)
=o(\sqrt n).
\label{eqn:broad-cost}
\end{align}

For an   $x^n\in\calE_n$ with $\|x^n\|^2\le B$, define its expected log-likelihood by
\begin{align}
\mu_{n,B}(x^n)
:=\mathbb E[\Lambda_B(x^n,Z_{\calE}^n)]=
r_n\rvC(B/r_n)
+\frac{\|x^n\|^2-B}{2(1+B/r_n)}.
\end{align}
Lemma~\ref{lem:broad_gauss} gives, uniformly over all $x^n\in\calE_n $ such that $\|x^n\|^2\le B$,
\begin{align}
\mu_{n,B}(x^n)\le \Upsilon_n(B),
\quad\mbox{and}\quad
\mathrm{Var}(\Lambda_B(x^n,Z_{\calE}^n))\le\frac{r_n}{2}.
\end{align}

Fix $\eta>0$. Since $\Upsilon_n(B)=o(\sqrt n)$, we have
$\eta\sqrt n-\Upsilon_n(B)>0$ for all sufficiently large $n$. Moreover,
because $\eta\sqrt n>0$, $ 
\{\Lambda_B(x^n,Z_{\calE}^n)^+>\eta\sqrt n\}
=
\{\Lambda_B(x^n,Z_{\calE}^n)>\eta\sqrt n\}.
$ 
Therefore, Chebyshev's inequality gives
\begin{align}
&\sup_{x^n\in\calE_n: \|x^n\|^2\le B}
\Pr\left(
\Lambda_B(x^n,Z_{\calE}^n)^+>\eta\sqrt n
\right)
\nonumber\\
&\quad\le
\sup_{x^n\in\calE_n: \|x^n\|^2\le B}
\frac{\mathrm{Var}(\Lambda_B(x^n,Z_{\calE}^n))}
     {(\eta\sqrt n-\mu_{n,B}(x^n))^2}
\nonumber\\
&\quad\le
\frac{r_n/2}{(\eta\sqrt n-\Upsilon_n(B))^2}
=O(n^{-5/8}) \to0.
\end{align}
This proves the assertion for both $B=b_n$ and $B=4b_n$.
\end{proof}

\end{document}
